\section*{Erd\H{o}s Problem 1124: the matching mechanism behind circle squaring}
\subsection*{Statement and scope}
Here \emph{circle} means the filled disk: a circular boundary has area zero and is not the intended object. Can a disk and a square of equal positive area be partitioned into finitely many pairwise congruent corresponding pieces? Laczkovich proved yes, using translations alone. No regularity of the pieces is required in this statement.
We prove the precise bounded-matching reduction, a necessary orbit-discrepancy estimate, and why finite convex pieces cannot suffice. We do not establish the uniform matching inequalities for a disk and a square. Their verification is the central missing step, so the resolution remains unresolved in this attempt, although the theorem is known.

\subsection*{A sufficient finite-neighbourhood condition}
Choose vectors $v_1,\ldots,v_d\in\mathbb R^2$ with no nonzero integer relation. Their integer span $\Gamma$ acts freely by translations. Fix $m\ge1$ and put
\[
 T=\left\{\sum_{j=1}^d n_jv_j:|n_j|\le m,\ n_j\in\mathbb Z\right\}.
\]
For every orbit $O=x+\Gamma$, consider the bipartite graph whose two vertex classes are labelled copies of $A\cap O$ and $B\cap O$, joining $a$ to $b$ when $b-a\in T$. Suppose that every finite subset of either class has at least as many neighbours in the other class.

Enumerate the countable first class. Finite Hall's marriage theorem supplies a matching for its first $N$ vertices for every $N$. Because every vertex has finitely many neighbours, the tree of these partial matchings is finitely branching with nodes at every depth. K\"onig's infinity lemma gives a matching saturating the first class. Applying the same argument in the other direction gives a matching saturating the second class. The usual Cantor--Bernstein construction combines the two injections into a bijection which uses an edge of one of these matchings at each matched pair. For finite classes use finite Hall directly. This proves the existence of a perfect matching.

Choose one such matching in each orbit, using choice. We obtain a bijection $F:A\to B$ with $F(a)-a\in T$. The sets
\[
 A_t=\{a\in A:F(a)=a+t\},\qquad t\in T,
\]
partition $A$, and the translated pieces $A_t+t$ partition $B$. There are at most $(2m+1)^d$ nonempty pieces. Conversely any translation equidecomposition using vectors in $T$ supplies such a matching and both finite-neighbourhood inequalities.

\subsection*{Why discrepancy enters}
Fix an orbit and identify it with $\mathbb Z^d$. If a matching has coordinate displacement at most $m$, then for $Q_L=\{-L,\ldots,L\}^d$,
\[
 |A\cap Q_{L-m}|\le|B\cap Q_L|,\qquad
 |B\cap Q_{L-m}|\le|A\cap Q_L|.
\]
For $L\ge m$ this gives
\[
 \bigl||A\cap Q_L|-|B\cap Q_L|\bigr|
 \le|Q_L\setminus Q_{L-m}|=O_{d,m}(L^{d-1}).
\]
Equal planar areas alone imply no such uniform statement on every countable orbit. Also a bound for cubes is not automatically the required Hall inequality for arbitrary finite sets. Both issues require proof in an existence argument.

\subsection*{Why finite convex pieces are impossible}
Suppose a disk could be equidecomposed with a square using finitely many convex pieces, allowing even rotations and reflections. Replace each piece by its closure. These compact convex sets cover the closed disk and have pairwise disjoint interiors; their congruent images similarly cover the square with pairwise disjoint interiors. Lower-dimensional pieces meet the circular boundary in at most two points each.

The boundary circle is covered by finitely many closed intersections with the pieces. The finite closed-cover form of Baire's theorem gives one piece whose closure contains a nonempty circular arc. Under its isometry, this arc lies in the square. A smaller open subarc lies in the square's interior, because a circle intersects the four boundary lines in only finitely many points.
Every point of this subarc belongs to another closed piece as well: it is a boundary point of the first piece, so approach it by points outside that piece but inside the square and pass to a subsequence belonging to one other piece. Another finite closed-cover argument shows that one other piece contains a nonempty subarc. Both convex pieces contain the convex hull of this subarc, which has nonempty planar interior, contradicting disjoint interiors.

Thus a finite decomposition necessarily uses some nonconvex pieces. This restriction on a simple class of pieces is compatible with the arbitrary pieces permitted by the problem.

\subsection*{Assessment and sources}
The matching criterion, the boundary discrepancy necessity and the convex-piece obstruction are proved. What is missing is a choice of translation vectors and uniform finite Hall inequalities for a disk and square. We do not replace that step by quoting the target theorem.
Sources: \url{https://www.erdosproblems.com/1124}; M. Laczkovich, \emph{Equidecomposability and discrepancy; a solution of Tarski's circle-squaring problem}, \url{https://doi.org/10.1515/crll.1990.404.77}; L. Grabowski, A. M\'ath\'e and O. Pikhurko, \emph{Measurable circle squaring}, Section 2 and Theorem 2.1, \url{https://annals.math.princeton.edu/wp-content/uploads/annals-v185-n2-p06-p.pdf}. The latter specifies the stronger summable uniform-discrepancy hypothesis used by the known solution. No novelty is claimed.
\subsection*{COMPLETION ESTIMATE}
\textbf{COMPLETION: 45\%}
